% Auto-rendered from open_questions.json; keep in sync.
\begin{enumerate}
  \item \textbf{Do the 2D PEPS chapters actually reproduce on a real problem at accessible bond dimension?}\\
        \emph{Basis.} This replication only exercised the 1D/MPS half of the notes; PEPS 2D benchmarks (Chap.~4) were left unrun because
        a laptop cannot fairly reach the $D\sim 7\text{--}10$ regime needed. Whether the notes' pedagogical PEPS pipeline actually recovers,
        e.g., QMC-level Heisenberg energies at moderate $D$, is untested and load-bearing.\\
        \emph{Next steps.} \texttt{quimb.tensor} PEPS simple-update on 2D Heisenberg, $L{=}6$, $D\in\{4..8\}$; compare per-site energy to
        published Sandvik QMC; free compute = uicgpu 8$\times$A100 (small-batch free lane). Success: $|E/N - E_\text{QMC}/N| < 5\times 10^{-3}$ at $D{=}8$.

  \item \textbf{Are the contraction-complexity scaling claims empirically achievable on modern hardware?}\\
        \emph{Basis.} We accepted the paper's asymptotic scaling (e.g.\ PEPS boundary-MPS $O(\chi^{10})$, MERA $O(\chi^9)$) without a single
        measured timing sweep. Memory locality / mixed-precision / contraction-order heuristics can shift the effective exponent on GPUs.\\
        \emph{Next steps.} Timing sweep: fixed lattice, $\chi\in\{8,16,24,32,48,64\}$, on (a) CPU \texttt{quimb}, (b) GPU
        \texttt{quimb}+\texttt{cotengra}+\texttt{jax}, (c) A100 direct. Fit $\log t$ vs $\log \chi$; compare slope to the notes' exponent.
        Free compute = uicgpu A100 idle. Deliverable: exponent $\pm$ CI per backend.

  \item \textbf{Do the canonicalization / optimal-truncation guarantees carry to fermionic TNs (swap gates or Grassmann tensors)?}\\
        \emph{Basis.} The paper touches fermionic TNs only briefly. Whether the C3 optimal-truncation identity
        $\text{err} = \sum_{k>\chi}\sigma_k^2$ still holds under fermion swap gates is not obviously true and was not tested here.\\
        \emph{Next steps.} Fermionic Hubbard 1D at $U{=}4$, half-filling; fMPS via \texttt{block2} or \texttt{quimb}'s fermionic backend;
        canonical-form check and truncation-error benchmark identical to C3. Free compute = laptop or uicgpu. Deliverable: yes/no on the C3 identity in the fermionic case.

  \item \textbf{Does the paper's clean-Ising entanglement pedagogy survive under disorder?}\\
        \emph{Basis.} The paper's C2 pedagogy is illustrated on the clean critical Ising chain. Whether the same MPS + slope-fit
        workflow still yields a meaningful "$c$" for a disordered chain (whose ground state flows to a Refael--Moore infinite-randomness
        fixed point) is a natural failure test of the notes' teaching.\\
        \emph{Next steps.} Redo the entropy experiment with $h_i \sim \mathcal{U}(0.5, 1.5)$, average over $\sim 50$ disorder samples at $N=64, 128$;
        compare fitted slope to $c_\text{eff} = (\ln 2)^2/6 \approx 0.080$. Free compute = cherryrd CPU / uicgpu.
        Success = clean separation between clean-chain and disorder-averaged slope.

  \item \textbf{Do the notes' DMRG results compose with modern autodiff-TN optimization (\texttt{jax}/\texttt{pytorch})?}\\
        \emph{Basis.} The notes predate the autodiff-TN wave. A modern reader will likely use a jax-based variational MPS energy-minimizer
        instead of a sweep-based DMRG; the paper does not certify that the two agree.\\
        \emph{Next steps.} Same TFIM $N{=}32$, $\chi{=}16$ problem; (a) \texttt{quimb} DMRG (our exp1), (b) \texttt{quimb}+\texttt{jax}
        autodiff MPS optimizer. Compare final energy and iteration count. Free compute = uicgpu A100. Deliverable: $|E_\text{DMRG} - E_\text{autodiff}|$
        and iteration-count ratio.
\end{enumerate}
